% BHU PYQ -- Transcribed & Corrected
% Paper: BOTMJ-43 / BOTMJ43: Gymnosperms and Paleobotany
% Exam: B.Sc. (Hons) Semester IV Examination, 2025-26 (NEP)
% Category: Major Core
% Department: Department of Botany, Institute of Science / Faculty of Science, Banaras Hindu University

\documentclass[11pt,a4paper]{article}
\usepackage[a4paper,top=2.5cm,bottom=2.5cm,left=2.8cm,right=2.8cm,headheight=14pt]{geometry}
\usepackage{amsmath,amssymb}
\usepackage[T1]{fontenc}
\usepackage[utf8]{inputenc}
\usepackage{lmodern,microtype,fancyhdr,enumitem,hyperref,lastpage,parskip}

\hypersetup{
    colorlinks=true,
    urlcolor=black,
    linkcolor=black,
    pdftitle={BOTMJ43: Gymnosperms and Paleobotany -- BHU B.Sc. Sem IV 2025-26 (NEP)}
}

\pagestyle{fancy}
\fancyhf{}
\renewcommand{\headrulewidth}{0.4pt}
\renewcommand{\footrulewidth}{0.4pt}
\fancyhead[L]{\small \textit{Banaras Hindu University}}
\fancyhead[R]{\small \textit{BOTMJ-43: Gymnosperms and Paleobotany (Sem IV 2025-26)}}
\fancyfoot[C]{\small Page \thepage\ of \pageref{LastPage}}

\newlist{parts}{enumerate}{2}
\setlist[parts,1]{label=(\alph*),leftmargin=*,itemsep=4pt,topsep=2pt}

\newcommand{\pts}[1]{\hfill \textit{[#1]}}

\begin{document}

\begin{center}
  {\large\bfseries BANARAS HINDU UNIVERSITY}\\[2pt]
  {\normalsize B.Sc. (Hons) Semester IV Examination, 2025-26 (NEP)}\\[6pt]
  \rule{0.8\linewidth}{0.5pt}\\[6pt]
  {\large\bfseries BOTANY}\\[4pt]
  {\large\bfseries Paper No. BOTMJ43 : Gymnosperms and Paleobotany}\\[2pt]
  {\small Ref. No.: Conf/Even/N/2025-26/20}\\[4pt]
  {\normalsize Time: $2\frac{1}{2}$ hours \hfill Full Marks: 70}
\end{center}

\rule{\linewidth}{0.4pt}

\smallskip

\noindent \textit{(Write your Roll No. at the top immediately on the receipt of this question paper)}\\[2pt]
\noindent \textit{The figures in the right-hand margin indicate marks.}\\[1pt]
\noindent \textit{Answer \textbf{five} questions including Question 1 which is compulsory.}

\bigskip

\noindent \textbf{Question 1.} Answer the following questions in not more than 100 words each: \pts{7 $\times$ 2 = 14}
\begin{parts}
  \item Why are the leaves of \textit{Cycas} described as pinnately compound but fern-like in appearance? \pts{2}
  \item Describe the salient features of Progymnosperm. \pts{2}
  \item Why are the leaves of \textit{Ephedra} reduced and what is the role of its green, jointed stems? \pts{2}
  \item Give an account of the affinities of Pteridospermales with Cycadophyta and ferns. \pts{2}
  \item Why are the seeds of \textit{Pinus} winged and how are these wings derived? \pts{2}
  \item Describe the anatomical features of \textit{Lyginopteris}. \pts{2}
  \item What is the significance of the presence of resin canals for survival in \textit{Pinus}? \pts{2}
\end{parts}

\bigskip
\rule{\linewidth}{0.2pt}
\bigskip

\noindent \textbf{Question 2.} Describe the geological time scale, outlining its major divisions, key events in each era and its significance in understanding the evolution of plant life. \pts{14}

\medskip\rule{\linewidth}{0.2pt}\medskip

\noindent \textbf{Question 3.} Write short notes on \textbf{any two} of the following: \pts{2 $\times$ 7 = 14}
\begin{parts}
  \item Reproductive parts of \textit{Lyginopteris} \pts{7}
  \item Pteridospermales \pts{7}
  \item Classification of gymnosperms \pts{7}
\end{parts}

\medskip\rule{\linewidth}{0.2pt}\medskip

\noindent \textbf{Question 4.} Give a comprehensive account of the types of fossils and their importance, with suitable diagrams. \pts{14}

\medskip\rule{\linewidth}{0.2pt}\medskip

\noindent \textbf{Question 5.} Write short notes on \textbf{any two} of the following: \pts{2 $\times$ 7 = 14}
\begin{parts}
  \item Reproduction in \textit{Pinus} \pts{7}
  \item Difference between young and old stems of \textit{Ephedra} \pts{7}
  \item Structure and adaptations of \textit{Pinus} leaf \pts{7}
\end{parts}

\medskip\rule{\linewidth}{0.2pt}\medskip

\noindent \textbf{Question 6.} Describe the anatomy and reproduction of \textit{Cycas}, highlighting its distinctive features and evolutionary significance. \pts{14}

\medskip\rule{\linewidth}{0.2pt}\medskip

\noindent \textbf{Question 7.} Write short notes on \textbf{any two} of the following: \pts{2 $\times$ 7 = 14}
\begin{parts}
  \item Major geological events in Earth's history \pts{7}
  \item \textit{Glossopteris} leaf \pts{7}
  \item Poikilohydry and homoiohydry \pts{7}
\end{parts}

\vfill
\rule{\linewidth}{0.4pt}

\begin{center}\small
\href{https://bhu-exam.vercel.app/}{Click here to visit our website}\\[4pt]
\href{https://me-aryan.vercel.app/}{Click here to visit our contributor}
\end{center}

\end{document}
